\section{Systematic Evaluation} \label{sec:eval}

Before we start evaluating the optimal object avoidance motion generator, a test setup is required to evaluate the performance of our solution.
In \cref{sec:introduction:comparability} we showed the inconsistency of test setups for a selection of state of the art controllers.

\subsection{Benchmark Definition} \label{sec:eval:setup}
These encapsulate a full pipeline, spanning object detection, tracking, motion generation and execution.
An isolated evaluation of the most important component, the motion generation, is therefore difficult to achieve.
\\
Systems like {PAC-MAN} \cite{yang_2026} are tested in real world or physics simulations.
Gazebo \cite{koenig_2004} or \acrfull{mujoco} \cite{todorov_2012} are two examples for such physical simulations.
A motion generation focused analysis of complete black-box controllers therefore always includes simulation errors.
As these physics simulations are approximations of the real-world, they model it in different ways and introduce errors for different use cases.
A great showcase of these differences are the contact formulations.
Gazebo allows for multiple physics engines, but comes by default with \acrfull{dart} \cite{lee_2018}.
\acrshort{dart} models contacts and collisions using an implicit time-stepping, velocity-based \acrfull{lcp} formulation.
In contrast \acrshort{mujoco} employs a non-complementary method, which turns it into a convex optimization problem.
So Gazebo has a rigid contact model, while \acrshort{mujoco} uses spring-damper systems as contact models \cite{lee_2018} \cite{todorov_2012}.
These fundamental differences remind us again that there is no such thing as a ground-truth in physics simulation and we should avoid evaluating with a single simulation.
As our work focuses on the collision avoidance motion generation we cut out all other components and evaluate solely if the motion generator finds a suitable joint trajectory.
\\
The resulting motions can later be executed with a suitable whole-body controller in a physics simulation or the real-world.
However, this would also include the errors of the chosen whole-body controller and possibly simulation.
\par
Our evaluation uses Karl Hupplers five characteristics of a good benchmark to derive properties our collision avoidance test setup has to implement.
From those we formulate requirements towards the start velocity, position and impact surface.
\\
These lead to the problem of uniformly sampling an ellipsoid, which is our selected impact surface.
The next section describes the impact location sampling process. (\cref{sec:eval:sampling})
With the impact locations in place, we present the full algorithm, including the calculation of object velocities and start locations in \cref{sec:eval:algo}.
\\
We conclude with an extension towards spline-based trajectories in \cref{sec:eval:extension}.

Our benchmark consists of multiple independent runs, where a single object approaches the robot.
The object is on a defined trajectory, targeted towards an impact location on the robot.
For our work we stick to linear object trajectories with constant velocity.
\\
A more complex definition of trajectories is also possible, however this would require small adjustments to the start position calculations.
An extension towards spline-based trajectories is described in \cref{sec:eval:extension}.
% \par
\paragraph{Required Benchmark Properties} \label{par:eval:setup:properties}
First we are deriving the necessary properties of a suitable benchmark setup.
"The Art of Building a Good Benchmark" \cite{huppler_2009} from Karl Huppler, mentions five characteristics: \textit{relevance}, \textit{repeatability}, \textit{fairness}, \textit{verifiability}, \textit{economicalness}, which need to be balanced in order to build a good benchmark.
While Hupplers characteristics are formulated for computer systems, we can use them to identify properties for our benchmark.
\\
A \textit{relevant} benchmark for motion generation of collision avoidance should represent an amount of successfully avoided objects, as this is "something important" to the user \cite{huppler_2009}.
One of the common metrics in this field is success rate \cite{yang_2026,zhang_2026}.
\\
To produce a success rate that truthfully reflects the strength and weaknesses of the motion generator for a specific robot a \textit{verifiable} setup is required.
In our context, \textit{verifiable} boils down to a testing method, which tests a range of collision avoidance scenarios, where the object approaches the robot from vastly different angles and covers various impact locations.
% Imagine a humanoid in its default upright standing pose.
% The method has to test impacts on the whole upper body surface with vastly differing impact angles.
\\
\textit{Repeatable} is in our case straightforward.
The benchmark should always result in the same success rate, when the generator itself is deterministic.
For \textit{economicalness} the benchmark has to be correct for any number of runs.
During ongoing development a small amount of individual runs is likely to be used, while for thorough in-depth evaluations a larger scale is required.
In general this is the weakest requirement, as long as it scales reasonably with the motion generator and does not introduce bias to the success rate.
\\
The last property is derived from \textit{fairness}.
Each evasion motion generator has to be able to perform on an even playing field.
This is another argument, why we are evaluating the existence of a found motion instead of executing it in a physics simulation.
Our evaluation consequently relies on the motion generators to implement a feasibility constraint and report the nonexistence avoidance motions.
\par
In summary our benchmark has to use runs, which are equally important, so that the success rate is \textit{verifiable}, while providing \textit{fairness}.
It should also be able to handle high and low amounts of test runs to accommodate for \textit{economicalness}.

\paragraph{Start Velocity, Position and Impact Surface} \label{par:eval:setup:velPosImpact}
The user of the benchmark has to select the impact surface $\mathcal I$, amount of objects $n_p$, time of flight $T_f$ and distance $d$ from start point to impact surface.
Start velocities $\bm v_i$ and impact locations are than calculated for each run $i$.
\par
To construct a \textit{verifiable} and \textit{relevant} benchmark from multiple single collision avoidance checks, all of them have to be equally important.
As a result every run needs to test a separate ability of the robot motion generator combination.
If runs are hitting similar impact locations or meet in a single focal point, the benchmark is weighting these areas stronger, as every run is equally weighted into the success rate.
\\
The solution is to sample uniformly spaced locations $\bm p_i^\mathcal I$ on the impact surface $\mathcal I$.
Normals $\bm n_i$ at the impact locations than define the impact direction.
In order to avoid creating focal points for the object trajectories an ellipsoid can be used.
Ellipsoids normal vectors meet in a focal cone instead of a single focal point \cite{staude_1886}, which concludes in an even penetration of the volume.
\par
Our use case is morpheus standing evasion, which results in an impact surface that is fitted around the torso.
It is however, also possible, to place the impact surface on other parts of the robot to evaluate collision avoidance abilities.
An example would be an impact surface at the lower leg for near ground level threat scenarios.
\\
\begin{figure}
  \centering
  \includegraphics[width=0.5\textwidth]{figures/05_benchmark/impactArea/cut-off.pdf}
  \caption{Cut-offs $\underline z$ and $\overline z$ in the xy-plane remove steep angles and enable the benchmark to focus on primarily horizontally approaching objects.}
  \label{fig:eval:setup:velPosImpact:cuf-off}
\end{figure}
As \cref{fig:eval:setup:velPosImpact:cuf-off} display, we define two vertical cut-offs $\overline z$ and $\underline z$, which are measured from the center line of the ellipsoid.
These enable the user to remove steep angles from the benchmark execution and puts the focus on horizontally approaching threats.
\\
The next section describes the impact location sampling on top of our impact surface.
Abstractly speaking, we are therefore formulating a method to uniformly sample $n_p$ points $\bm p_i^\mathcal I$ on an ellipsoid ring $\mathcal I$, which is characterized by the ellipsoid radius $\bm r_\mathcal{I}$ and cut-offs $\underline z$, $\overline z$.
The following sections present our uniform ellipsoid sampling algorithm.

\subsubsection{Ellipsoid Sampling} \label{sec:eval:sampling}
\begin{figure}
  \centering
  \includegraphics[width=0.7\textwidth]{figures/05_benchmark/gonzalez/comparison.pdf}
  \caption{Comparison between Latitude-longitude lattices (top) and Fibonacci lattice (bottom) from Gonz{'a}lez 2010 \cite{gonzalez_2010}.}
  \label{fig:eval:sampling:sampling:latticeComp}
\end{figure}
\Cref{fig:eval:sampling:sampling:latticeComp} shows a comparison between classical Latitude-longitude lattices and Fibonacci lattices.
While Latitude-longitude lattices produce nonuniform positions, which are clustered at the poles, Fibonacci lattices are a great option for uniformly selecting samples on a sphere \cite{gonzalez_2010}.
% Gonz{'a}lez points in his 2010 research \cite{gonzalez_2010} out that Fibonacci lattices are a great option for uniform sampling on spheres.
However, Fibonacci lattices samples lose their equal distancing when transformed to ellipsoids \cite{marples_2024}.
\\
In general, rejection sampling methods, like gradient vector rejection sampling or area element rejection sampling, are advised for ellipsoids as they offer great solutions for high sample counts \cite{marples_2024,williamson_1987,melfi_2004}.
As rejection sampling is driven by random point selection and rejection of none suitable samples, research like Marples et al. \cite{marples_2024} mostly compares these sampling-based methods on high sample counts (e.g., $10^{8}$).
For low sample sizes, those methods show overpowering noise and therefore are not sufficient for our sampling requirements.
\\
\begin{figure}
  \centering
  \begin{subfigure}[b]{0.3\textwidth}
    \includegraphics[draft,width=\textwidth,height=5cm]{figures/05_benchmark/ZZ_results/compare/rejection_1e2.pdf}
    \caption{Rejection Sampling ($10^2$)}
    \label{fig:eval:sampling:compare:rejection:low}
  \end{subfigure}
  \hspace{4em}
  \begin{subfigure}[b]{0.3\textwidth}
    \includegraphics[draft,width=\textwidth,height=5cm]{figures/05_benchmark/ZZ_results/compare/rejection_1e8.pdf}
    \caption{Rejection Sampling ($10^8$)}
    \label{fig:eval:sampling:compare:rejection:high}
  \end{subfigure}
  \\
  \vspace{2em}
  \begin{subfigure}[b]{0.3\textwidth}
    \includegraphics[draft,width=\textwidth,height=5cm]{figures/05_benchmark/ZZ_results/compare/fibonacci_1e2.pdf}
    \caption{Fibonacci Lattices ($10^2$)}
    \label{fig:eval:sampling:compare:fibonacci:low}
  \end{subfigure}
  \hspace{4em}
  \begin{subfigure}[b]{0.3\textwidth}
    \includegraphics[draft,width=\textwidth,height=5cm]{figures/05_benchmark/ZZ_results/compare/fibonacci_1e8.pdf}
    \caption{Fibonacci Lattices ($10^8$)}
    \label{fig:eval:sampling:compare:fibonacci:high}
  \end{subfigure}
  \caption{
  Comparison between gradient vector rejection and Fibonacci lattices with an ellipsoid reparametrization.
  We compute $10^2$ (left) and $10^8$ (right) samples.
  It is clearly visible, that the sampling in (a) is dominated by noise.
  }
  \label{fig:eval:sampling:compare}
\end{figure}
This can be seen in \cref{fig:eval:sampling:compare:rejection:low,fig:eval:sampling:compare:rejection:high} were, we display the uniformity of gradient rejection sampling for $10^2$ and $10^8$ samples.
We can see the noise vanishing for high sample sizes, while dominating the resulting uniformity in \cref{fig:eval:sampling:compare:rejection:low}.
In comparison sampling with Fibonacci lattices on a unit sphere and than transforming to an ellipsoid, achieves good results for low sample sizes, while introducing non uniformity for high samples (\cref{fig:eval:sampling:compare:fibonacci:low,fig:eval:sampling:compare:fibonacci:high}).
To quantify uniformity we used the in Marples et al. 2024 \cite{marples_2024} defined points per segment metric.
\todo[inline]{describe area weighted \acrshort{rsd}}
\begin{figure}
  \centering
  \includegraphics[width=0.4\textwidth]{figures/05_benchmark/marples/segments.pdf}
  \caption{
  This figure from Marples et al. 2024 \cite{marples_2024} displays their segment definition for setting a quantifiable uniformity metric.
}
  \label{fig:eval:sampling:uniformitySegments}
\end{figure}
Segments are defined like shown in \cref{fig:eval:sampling:uniformitySegments}, which was taken from Marples et al. 2024 \cite{marples_2024}.
For the comparison and noise analysis in \cref{fig:eval:sampling:compare} we used $\Delta  \theta = \SI{1}{\deg}$ and $\Delta  \Phi = \SI{1}{\deg}$.
\par
We propose a combination of Fibonacci Lattices \cite{gonzalez_2010} with the optimization strategy of the well-known Thomson Problem \cite{thomson_1904}, to achieve a uniform sampling technique, that is also valid for small sample sizes.
First our sampling places points on a sphere with the Fibonacci lattice strategy, which is described in \cref{sec:eval:sampling:fibonacci}.
The sphere is than stretched to an ellipsoid.
After transformation the points on the ellipsoid are not uniformly spaced anymore.
However, as they are still roughly placed in the correct spot, we use them as starting points for our optimization (\cref{sec:eval:sampling:thomson}).

\paragraph{Fibonacci Lattices} \label{sec:eval:sampling:fibonacci}
\begin{figure}
  \centering
  \includegraphics[width=0.7\textwidth]{figures/05_benchmark/gonzalez/algorithm.pdf}
  \caption{
  Fibonacci lattice with 21 samples on a Fermal spiral.
  Red: golden angle; Blue: complementary golden angle. \cite{gonzalez_2010}
  }
  \label{fig:eval:sampling:fibonacci:algoSpiral}
\end{figure}
Fibonacci Lattices are constructed by leveraging the well-known golden angle to pick points on a spiral.
Gonz{'a}lez picked the Fermat spiral \cite{vogel_1979} to have an equal angle to area turn, while provindg an offset to the poles.
This spiral, combined with the golden angle, produces a sampling that spreads the points across the whole sphere, while resulting in isotropic spacing \cite{gonzalez_2010}.
\\
\Cref{fig:eval:sampling:fibonacci:algoSpiral} displays the Fermat spiral.
The figure shows two spirals, the red one was created with the golden angle, while a complementary is colored blue.
Both spirals produce the same sampling as \cref{fig:eval:sampling:fibonacci:algoSpiral} shows.
For a much deeper analysis, we strongly recommend Gonz{'a}lez 2010 work; it presents in great detail the construction and uniformity of Fibonacci Lattices.
\\
Another advantage of these lattices, compared with rejection sampling methods, is their deterministic nature.
The construction includes no randomness and produces the same point positions for every run; therefore, no pseudo-random seeding is required at this step.
% We strive for a uniform point placement on an ellipsoid ring with great reproducibility and accuracy for low sample sizes.
To move from this unit sphere to an ellipsoid, we stretch the Fibonacci Lattice generated samples.

\paragraph{Inter-point Distance Optimization} \label{sec:eval:sampling:thomson}
During the process of reparameterizing the sphere, we changed the spacing between our samples (\cref{fig:eval:sampling:overview:transform}).
They are no longer uniformly placed on the ellipsoid.
For a typical humanoid upper body, the hull ellipsoid is stretched upwards and compressed around the frontal plane.
Because of the cut-offs in \cref{fig:eval:setup:velPosImpact:cuf-off}, we ignore the high error regions at the top and bottom.
The horizontal compression, however, introduces clumping of points at the sides of the robot, while fewer points are found at the front and back.
\\
We trim down both errors by applying optimization, which is also used to solve the Thomson problem \cite{thomson_1904} or also known as electrons on a sphere problem.
The solution in Dolan et al. 2004 \cite{dolan_2004} minimizes the electrostatic potential energy, which results in an even spread of electrons across the unit sphere.
In our case it maximizes the accumulated inter-point distances for impact locations $\bm p_i^\mathcal I$ between the upper $\overline z$ and lower ring cut-off $\underline z$ on our ellipsoid impact surface.
\Cref{eq:eval:sampling:thomson} is the adapted formulation for the point count $n_p$.
\begin{subequations}
\label{eq:eval:sampling:thomson}
  \begin{align}
    \min_{\bm p} \sum^{n_p-1}_{i=1} \sum^{n_p}_{j=i+1} & \frac{1}{\|\bm p_i^{\mathcal I}- \bm p_j^{\mathcal I}\|_2} & \label{eq:eval:thomson:obj} \\
  \text{s.t.} \qquad \left|\left|\frac{\bm p_i^\mathcal I}{\bm r}\right|\right|_2^2 & = 1 & \forall i \in \left[ 1, n_p \right] \label{eq:eval:thomson:ellipsoid} \\
\overline z \leq \bm e_3^T \cdot p_i^{\mathcal I} & \leq \underline z &  \forall i \in \left[ 1, n_p \right] \label{eq:eval:thomson:cut}
  \end{align}
\end{subequations}
We adjusted the formulation from Dolan et al. 2004 \cite{dolan_2004} to \cref{eq:eval:sampling:thomson} with an additional cut-off constraint (\cref{eq:eval:thomson:cut}) and reparametrization (\cref{eq:eval:thomson:ellipsoid}). \\
For our sampling method, we set the stretched Fibonacci lattice points as initial guesses for the problem, as they are already sufficiently close to the solution, resulting in a refined positioning of the original points.
In order to produce a reproducible sampling, the optimization algorithm has a fixed seed of zero.
NLopt \cite{NLopt}, with its implementation of the \acrshort{slsqp} \cite{kraft_1994} solver, was chosen for solving \cref{eq:eval:sampling:thomson}.

\paragraph{Resulting Sampling Algorithm} \label{sec:eval:sampling:algo}
We proposed an ellipsoid sampling algorithm to produces a uniform spread of points $\bm p_i^\mathcal I$ across the chosen impact ellipsoid surface $\mathcal I$.
\\
\begin{figure}
  \centering
  \begin{subfigure}[b]{0.3\textwidth}
    \includegraphics[draft,width=\textwidth,height=5cm]{figures/05_benchmark/ZZ_results/sampling/fibonacci.pdf}
    \caption{Fibonacci Lattices}
    \label{fig:eval:sampling:overview:fibonacci}
  \end{subfigure}
  \hfill
  \begin{subfigure}[b]{0.3\textwidth}
    \includegraphics[draft,width=\textwidth,height=5cm]{figures/05_benchmark/ZZ_results/sampling/transform.pdf}
    \caption{Transform to Ellipsoid}
    \label{fig:eval:sampling:overview:transform}
  \end{subfigure}
  \hfill
  \begin{subfigure}[b]{0.3\textwidth}
    \includegraphics[draft,width=\textwidth,height=5cm]{figures/05_benchmark/ZZ_results/sampling/thomson.pdf}
    \caption{Optimized}
    \label{fig:eval:sampling:overview:thomson}
  \end{subfigure}
  \caption{
  To compute impact locations $\bm p_i^\mathcal I$ on the surface $\mathcal I$, we first sample on a unit sphere, using Fibonacci lattices (a).
  Afterwards we reparameterize the unit sphere with the ellipsoid radius $\bm r$ (b).
  The impact locations are not uniformly spread on the impact surface anymore.
  We apply an optimization strategy that is also used to solve the Thomson problem and restore the uniformity of the samples.
  }
  \label{fig:eval:sampling:overview}
\end{figure}
The algorithm starts of by sampling a unit sphere by utilizing Fibonacci lattices \cite{gonzalez_2010} (\cref{fig:eval:sampling:overview:fibonacci}).
Afterwards we transform the underlying sphere and sampled points to an ellipsoid (\cref{fig:eval:sampling:overview:transform}).
\\
To reinstate the introduced uniformity, we optimize the inter-point distance (\cref{fig:eval:sampling:overview:thomson}).
This last step is analog to solving the Thomson problem \cite{thomson_1904}, with a formulation like the one mentioned in Dolan et al. 2004 \cite{dolan_2004}.
We extended Dolans optimization with an ellipsoid constraint and cut-off limitations $\overline z$ and $\underline z$.
The full optimization can be found in \cref{eq:eval:sampling:thomson}.
For the implementation NLopt \cite{NLopt}, with its implementation of the \acrshort{slsqp} \cite{kraft_1994} solver, was chosen.
\\
Next up we present the complete benchmark algorithm of which the previous section is the core part.

\subsubsection{Algorithm Structure} \label{sec:eval:algo}
We derived an ellipsoid sampling method, which is able to select uniformly on low sample counts.
It is used to extract $n_p$ impact locations $\bm p_i^\mathcal I$ from the user defined surface $\mathcal I$.
\begin{figure}
  \centering
  \begin{subfigure}[b]{0.45\textwidth}
    \includegraphics[width=\textwidth]{figures/05_benchmark/impactArea/front.pdf}
    \caption{Front View (YZ)}
    \label{fig:eval:algoStructure:impactArea:front}
  \end{subfigure}
  \hfill
   \begin{subfigure}[b]{0.45\textwidth}
    \includegraphics[width=\textwidth]{figures/05_benchmark/impactArea/side.pdf}
    \caption{Side View (XZ)}
    \label{fig:eval:algoStructure:impactArea:side}
  \end{subfigure}
  \vspace{1em}
  \\
  \begin{subfigure}[b]{0.45\textwidth}
    \includegraphics[width=\textwidth]{figures/05_benchmark/impactArea/top.pdf}
    \caption{Top View (XY)}
    \label{fig:eval:algoStructure:impactArea:top}
  \end{subfigure}
  \hfill
 \begin{subfigure}[b]{0.5\textwidth}
    \includegraphics[width=\textwidth]{figures/05_benchmark/impactArea/velocity.pdf}
    \caption{Overview}
    \label{fig:eval:algoStructure:impactArea:overview}
  \end{subfigure}
  \caption{
    The red points shows the tightly fitted impact hull around the robot.
    Blue mark the outwards projected start positions for the approaching objects.
    The trajectory for each individual projectile from the starting location (blue) through the chosen impact point (red) is visualized in gray.
  }
  \label{fig:eval:algoStructure:impactArea}
\end{figure}
\Cref{fig:eval:algoStructure:impactArea} displays $500$ impact locations as red points around the blue colored collision hull of the robot.
The impact surface was chosen to cover the upper body and head area with a dense grid.
We cut the surface on the top and bottom to primarily test for horizontally approaching objects.
\par
From the sampled impact locations we are able to derive velocities $\bm v$ and $\bm p_i^0$.
In \cref{fig:eval:algoStructure:impactArea} object start positions $\bm p_i^0$ are marked as blue points.
The distance $d$ between impact location $\bm p_i^\mathcal I$ and $\bm p_i^0$ is set by the user to control the trajectory.
% start position p_0 based on d
Velocities are in the direction of surface normals $\bm n_i$.
% velocity magnitude based on d
\par
% motion generation and success rate calculation
\par
Together with the previously described benchmarking method, any humanoid collision avoidance motion generator can be tested and investigated for its limitations.

\subsection{Evaluation of the Optimal Motion Generator} \label{sec:eval:app}
\subsubsection{Results} \label{sec:eval:app:result}

\begin{figure}
  \centering
  \begin{subfigure}[b]{0.45\textwidth}
    \includegraphics[width=\textwidth]{figures/05_benchmark/ZZ_results/conv/top.pdf}
    \caption{Top View (XY)}
    \label{fig:eval:results:convTop}
  \end{subfigure}
  \hfill
  \begin{subfigure}[b]{0.45\textwidth}
    \includegraphics[width=\textwidth]{figures/05_benchmark/ZZ_results/conv/iso.pdf}
    \caption{Isometric View}
    \label{fig:eval:results:convIso}
  \end{subfigure}
  \vfill
  % \hfill
  \begin{subfigure}[b]{0.45\textwidth}
    \includegraphics[width=\textwidth]{figures/05_benchmark/ZZ_results/conv/front.pdf}
    \caption{Front View (YZ)}
    \label{fig:eval:results:convFront}
  \end{subfigure}
  \hfill
  \begin{subfigure}[b]{0.45\textwidth}
    \includegraphics[width=\textwidth]{figures/05_benchmark/ZZ_results/conv/side.pdf}
    \caption{Side View (XZ)}
    \label{fig:eval:results:convSide}
  \end{subfigure}
  \caption{orange: no convergence, gray lines are not converged runs}
  \label{fig:eval:results:conv}
\end{figure}
The benchmark results from 500 individual test runs shown earlier are visualized in \cref{fig:eval:results:conv}.
The figure displays the starting position of each object, as well as the impact locations of trajectories that could not be evaded.
Green marks a successfully generated avoidance trajectory, while orange shows occasions where the solver could not find an optimal motion.
These failed runs originate from different causes, ranging from significantly exceeding the typical iteration count to divergence in the \acrfull{qp} problem.
Overall, our optimal trajectory generator achieved an 81.2\% (406) success rate.

\begin{figure}
  \centering
  \includegraphics[width=\textwidth]{figures/05_benchmark/ZZ_results/plots/torqueJoint_original.pdf}
  \caption{
    Joint angles $\bm{q}$ and torques $\bm{\tau}$ for all successful dodging motions.
    We can clearly see that some of those trajectories require the robot to operate at its limits.
    Torques in general show a high usage of the available strength, indicating that the success rate is indeed limited by torque limits.
  }
  \label{fig:eval:results:limits}
\end{figure}

\begin{figure}
  \centering
  \includegraphics[width=0.7\textwidth]{figures/05_benchmark/ZZ_results/plots/distance.pdf}
  \caption{Distance of traveling object throughout the evasion motions over all successfully evaded object trajectories, with the minimal distance being clearly above the safety distance to the collision hull}
  \label{fig:eval:results:distance}
\end{figure}

Before we take a closer look at the non-converged cases, we visualize in \cref{fig:eval:results:distance} the distance between the robot and the projectile at each time step.
\Cref{fig:eval:results:distance} displays the sequence of distances for every individual successful object trajectory.
For better visualization, a convex hull around all sequences is drawn.
All generated motions come closest to the robot around the middle of the trajectory.
We also view the safety distance $d_\text{safe}$ in the form of a blue limit.
As is typical for optimal control, the solution found runs the robot as close as allowed to the limitations, sometimes even going down to the minimal allowed distance.
Another interesting insight is that for a minority of trajectories the robot is closest to the object before or after the mid time.
These are probably object trajectories, which are extremely steep.
For such evasion maneuvers, a leaning motion with long durations of close proximity to the object is the result.
It is clear that the optimal collision-avoidance motion generator goes right up to the set constraint.
\\
To better understand limitations imposed by either the trajectory generator or Unitrees H1-2, we investigate the impact locations and velocity vectors of the non-converged trajectories.
\Cref{fig:eval:results:conv} marks failed trajectory generation attempts as orange object starting points.
Most trajectories our controller struggles with come directly from the front or back and cross nearly the same area.
\\
To investigate potential hardware restrictions, we examine the motion-defining torque controls in \cref{fig:eval:results:limits}.
As in the distance plot \cref{fig:eval:results:distance}, we display controls for each converged motion.
A convex hull helps to visualize the extreme values of all controls.
The hardware torque limit for each joint is shown on the side of each plot, with minima and maxima marked for clarity.
We find that the robot is relying heavily on ankle and knee torque to stabilize and balance itself, as all reach their boundaries over an excessive time frame.
While hips and knees seem to be sufficiently powered, especially the ankle torques possibly limit the abilities of the robot to perform more demanding avoidance motions.
\\
A likely explanation lies in the general strategy we employ.
Most collision avoidance work on humanoids like \cite{yang_2026,kohlbrenner_2026} allows multiple capture steps.
Therefore, the robot does not need to apply high torques on its ankle as it steps away from the approaching object.
Our strategy, however, keeps the robot in a bipedal stance in order to avoid risking collisions with the environment, particularly other platforms possibly roaming the floor, requesting higher ankle and knee torques.
\\
Finally, we inspect the range of motion through \cref{fig:eval:results:limits}.
While the hip joints offer a huge range of motion, which is barely used, ankles and especially knee joints are driven towards their joint boundaries.
It is reasonable to assume that having the humanoid bend over could violate its dynamic stability constraint; therefore, the robot exploits squatting or upper body rotation to avoid collisions.
Crouching, in particular, is bound by the lower joint limit of the knees, as shown in \cref{fig:eval:results:limits}. \par
After we analyzed the generated motions, we found four distinct types, as well as one clear hardware limitation of the employed H1-2 for fixed-feet bipedal object evasion.
Our benchmark, however, provides a tool to make such hurdles visible.

\subsubsection{Motion Classes} \label{sec:eval:app:motionClasses}
\begin{figure}
  \centering
  \begin{subfigure}[b]{0.23\textwidth}
    \includegraphics[width=\textwidth]{figures/05_benchmark/ZZ_results/motions/crouch.png}
    \caption{crouching evasion}
    \label{fig:eval:results:motion:crouch}
  \end{subfigure}
  \hfill
  \begin{subfigure}[b]{0.23\textwidth}
    \includegraphics[width=\textwidth]{figures/05_benchmark/ZZ_results/motions/lean.png}
    \caption{leaning evasion}
    \label{fig:eval:results:motion:lean}
  \end{subfigure}
  % \vfill
  \hfill
  \begin{subfigure}[b]{0.23\textwidth}
    \includegraphics[width=\textwidth]{figures/05_benchmark/ZZ_results/motions/inverseNeo.png}
    \caption{reverse neo evasion}
    \label{fig:eval:results:motion:inverseNeo}
  \end{subfigure}
  \hfill
  \begin{subfigure}[b]{0.23\textwidth}
    \includegraphics[width=\textwidth]{figures/05_benchmark/ZZ_results/motions/neo.png}
    \caption{classic neo evastion}
    \label{fig:eval:results:motion:neo}
  \end{subfigure}
  \caption{
    We found, that our optimal object avoidance generator explored mainly four motion classes, for different kinds of approaching objects.
    The leaning evasion in (b) was found to be the most common one, while the reverse neo motion was leveraged in a small amount of object avoidance motions.
  }
  \label{fig:eval:results:motion}
\end{figure}
To start off, we explore the different strategies our optimal object avoidance controller found.
\Cref{fig:eval:results:motion} displays the four major evasion motions.
By far the most common motion is a leaning pose perturbation.
The robot mainly adjusts its upper body by tilting it away from the incoming object (\cref{fig:eval:results:motion:lean}).
Two interesting strategies result from low-flying objects, orthogonally targeted towards the robot's front or back (\cref{fig:eval:results:motion:inverseNeo,fig:eval:results:motion:neo}).
These object trajectories can be evaded by leaning back significantly, almost like the protagonist of the famous movie "The Matrix".
We therefore call these extremely dynamic strategies reverse and classic neo evasion.
Additionally, both leverage a small rotation around the torso joint to keep the robot's balance.
The last motion type is a crouching or squatting strategy.
The humanoid crouches underneath low horizontal flying targets, as seen in \cref{fig:eval:results:motion:crouch}.
For further visualization we provide a video with exemplary trajectories of the mentioned motion classes.
We suspect a weakness in the robot's ability to avoid front- and back-targeted objects, as those need to be evaded by employing difficult and highly dynamic strategies, like crouching, as well as the neo motions.
\\

